\section{References}

\begingroup\footnotesize

\subsection{Controlling liturgical books and present legislation}

\begin{itemize}
  \item \emph{Missale Romanum ex decreto Sacrosancti Oecumenici Concilii Vaticani II instauratum auctoritate Pauli PP.~VI promulgatum Ioannis Pauli PP.~II cura recognitum}, editio typica tertia, reimpressio emendata (Vatican City: Libreria Editrice Vaticana, 2008), especially \emph{Institutio Generalis Missalis Romani}, \emph{Ordo Missae}, Eucharistic Prayers I--IV, the prayers for Reconciliation, the four forms for Various Needs, and Appendix II\@. The Dicastery for Divine Worship's \href{https://www.cultodivino.va/en/formazione/pubblicazioni/libri-liturgici/missale/missale-romanum.html}{official bibliographic page} controls edition identity and subsequent variations.
  \item \emph{The Roman Missal}, Third Edition, English translation according to the use of the dioceses of the United States of America (Washington, DC: United States Conference of Catholic Bishops, 2011), identified but not reproduced. The \href{https://www.usccb.org/prayer-and-worship/the-mass/order-of-mass}{USCCB Order-of-Mass overview} and current liturgical-book notices control territorial scope.
  \item United States Conference of Catholic Bishops, \emph{General Instruction of the Roman Missal}, especially nos.~27--45 on the celebration's nature and elements; 46--54, Introductory Rites; 55--71, Word; 72--89, Eucharist; 90, conclusion; 91--111, ministries; 112--287, forms and ceremonial; 281--287, both species; 352--367, choices; and 368--399, Masses for Various Needs and adaptations. \href{https://www.usccb.org/prayer-and-worship/the-mass/general-instruction-of-the-roman-missal}{Current official U.S. text and chapter links}. Particular U.S. controls used include nos.~43, 48, 61, 87, 154, 160, 283, and 393. USCCB Committee on Divine Worship, \href{https://www.usccb.org/resources/newsletter-2021-02.pdf}{\emph{Newsletter} 57, no.~2 (February 2021), pp.~6--8}, records CDWDS confirmation Prot.~N.~228/20 and the U.S. decree of publication for the GIRM 54 emendment.
  \item \emph{Ordo Lectionum Missae}, editio typica altera (Vatican City: Libreria Editrice Vaticana, 1981), together with its \emph{Praenotanda}. The Dicastery's \href{https://www.cultodivino.va/en/formazione/pubblicazioni/libri-liturgici/missale/ordo-lectionum-missae.html}{official edition page} controls identification. The territorial books are the \emph{Lectionary for Mass for Use in the Dioceses of the United States of America}, second typical edition (1998/2001), with the \emph{Lectionary for Mass Supplement} (2017). USCCB, \href{https://www.usccb.org/news/2026/recap-us-bishops-plenary-assembly-orlando}{``Recap of U.S. Bishops' Plenary Assembly in Orlando'' (12 June 2026)}, records approval of elements for a future edition and the outstanding \latin{confirmatio} and \latin{recognitio}; it does not make that project operative.
  \item \emph{Ordo Cantus Missae}, editio typica altera (Vatican City, 1988); \emph{Graduale Romanum} (Solesmes, 1974); \emph{Graduale Simplex}, editio typica altera (Vatican City, 1988); and \emph{Caeremoniale Episcoporum}, editio typica (Vatican City, 1984), especially nos.~74 and 141. These control chant ordering and the episcopal branches actually invoked; they are not used to turn a local music program or pontifical ceremonial into the ordinary parish norm.
  \item Congregation for Divine Worship and the Discipline of the Sacraments, decree \emph{Paternas vices} (1 May 2013), Prot.~N.~215/11/L, inserting Saint Joseph into Eucharistic Prayers II--IV. \href{https://www.usccb.org/prayer-and-worship/the-mass/order-of-mass/liturgy-of-the-eucharist/saint-joseph-in-the-eucharistic-prayers}{Official U.S. decree and implementation record}; USCCB, \href{https://www.usccb.org/resources/newsletter-2013-07.pdf}{Committee on Divine Worship Newsletter, July 2013}, for mandatory U.S. use from 19 September 2013.
  \item Congregation for Divine Worship, \emph{Directory for Masses with Children} (1 November 1973); the separately approved \emph{Eucharistic Prayers for Masses with Children} for U.S. use, confirmed with Prot.~N.~692/11/L; and USCCB, \href{https://www.usccb.org/prayer-and-worship/the-mass/order-of-mass/liturgy-of-the-eucharist/eucharistic-prayer-for-masses-with-children}{current scope and restrictions}. The \href{https://www.usccb.org/resources/newsletter-2011-07.pdf}{July 2011}, \href{https://www.usccb.org/resources/newsletter-2011-08.pdf}{August 2011}, and \href{https://www.usccb.org/resources/newsletter-2022-04.pdf}{April 2022} USCCB newsletters control approval chronology, the separate 2011 supplement, and current-book status.
  \item Congregation for Divine Worship and the Discipline of the Sacraments, instruction \emph{Redemptionis Sacramentum} (25 March 2004), especially nos.~38--42, 51--54, 69--72, 88--96, 100--107, and 154--160. \href{https://www.vatican.va/roman_curia/congregations/ccdds/documents/rc_con_ccdds_doc_20040423_redemptionis-sacramentum_en.html}{Holy See English text}.
  \item Code of Canon Law (1983), especially canon 767 on the homily and canons 897--958 on the Eucharist, including ministers and recipients, participation, Eucharistic matter, offering, reservation, and the eucharistic fast. \href{https://www.vatican.va/archive/cod-iuris-canonici/eng/documents/cic_lib4-cann834-1253_en.html}{Holy See English text}. Canonical references state general law and do not adjudicate an individual's disposition.
\end{itemize}

\subsection{Conciliar, magisterial, and doctrinal controls}

\begin{itemize}
  \item Second Vatican Council, constitution \emph{Sacrosanctum Concilium} (4 December 1963), especially nos.~7, 14, 21--23, 26--35, and 47--58: Christ's presence, participation, ordered ministries, ritual reform, Word and Eucharist, homily, common prayer, Communion, concelebration, and vernacular. \href{https://www.vatican.va/archive/hist_councils/ii_vatican_council/documents/vat-ii_const_19631204_sacrosanctum-concilium_en.html}{Holy See English text}.
  \item Paul VI, apostolic constitution \emph{Missale Romanum} (3 April 1969), the reform's official account of the Roman Canon, three new Eucharistic Prayers, standardized institution narrative, memorial acclamation, richer Scripture, responsorial Psalm, restored common prayer, revised offering and fraction, and promulgation. \href{https://www.vatican.va/content/paul-vi/en/apost_constitutions/documents/hf_p-vi_apc_19690403_missale-romanum.html}{Holy See English text}.
  \item Second Vatican Council, dogmatic constitution \emph{Lumen gentium}, especially nos.~1, 10--11, 20--29, and 48--51, on the Church as sacrament, baptismal and ministerial priesthood, hierarchy and ministries, Eucharistic participation, and communion with the saints; dogmatic constitution \emph{Dei Verbum}, no.~21, on the Church's veneration of Scripture and the one table of Word and Christ's Body; decree \emph{Presbyterorum ordinis}, nos.~2 and 5. \href{https://www.vatican.va/archive/hist_councils/ii_vatican_council/documents/vat-ii_const_19641121_lumen-gentium_en.html}{Holy See text of \emph{Lumen gentium}}.
  \item Second Vatican Council, pastoral constitution \emph{Gaudium et spes} (7 December 1965), nos.~1, 4, and 11, on the Church's solidarity with human joy and anguish and Gospel-governed discernment of the signs of the times. \href{https://www.vatican.va/archive/hist_councils/ii_vatican_council/documents/vat-ii_const_19651207_gaudium-et-spes_en.html}{Holy See English text}.
  \item Council of Trent, Session XIII, decree on the Eucharist, especially chs.~1--8 and canons 1--4; Session XXI, doctrine on Communion under either species; Session XXII, doctrine and canons on the Sacrifice of the Mass, especially chs.~1--9. These control Real Presence, transubstantiation, sacramental sacrifice, the whole Christ under either species, and legitimate ecclesial discipline rather than an assertion that every sixteenth-century rubric was immutable.
  \item \emph{Catechism of the Catholic Church}, nos.~1066--1209 and 1322--1419, especially 1088, 1140--1144, 1345--1355, 1362--1381, 1391--1405, and 1408--1419. \href{https://www.vatican.va/archive/ENG0015/_INDEX.HTM}{Holy See English edition}.
  \item Pius XII, encyclical \emph{Mediator Dei} (20 November 1947), especially nos.~20--23, 47--64, 68--78, and 80--104, on the public worship of the Mystical Body, development, rejection of indiscriminate antiquarianism, Eucharistic sacrifice, and participation. \href{https://www.vatican.va/content/pius-xii/en/encyclicals/documents/hf_p-xii_enc_20111947_mediator-dei.html}{Holy See English text}.
  \item Paul VI, encyclical \emph{Mysterium Fidei} (3 September 1965), especially nos.~4--5, 24--27, 39--46, and 56--66, on Eucharistic mystery, sacrifice, transubstantiation, enduring Presence, worship, and unity. \href{https://www.vatican.va/content/paul-vi/en/encyclicals/documents/hf_p-vi_enc_03091965_mysterium.html}{Holy See English text}.
  \item John Paul II, encyclical \emph{Ecclesia de Eucharistia} (17 April 2003), especially nos.~11--20, 28--33, and 47--52, on sacrificial memorial, Presence, apostolic ministry, Communion, and liturgical form. \href{https://www.vatican.va/content/john-paul-ii/en/encyclicals/documents/hf_jp-ii_enc_20030417_eccl-de-euch.html}{Holy See English text}.
  \item Benedict XVI, apostolic exhortation \emph{Sacramentum Caritatis} (22 February 2007), especially nos.~6--15, 34--55, and 64: trinitarian gift, institution, sacrifice, participation, ars celebrandi, structure, dismissal and mission, and the method of mystagogy. \href{https://www.vatican.va/content/benedict-xvi/en/apost_exhortations/documents/hf_ben-xvi_exh_20070222_sacramentum-caritatis.html}{Holy See English text}.
  \item Francis, apostolic letter \emph{Desiderio desideravi} (29 June 2022), especially nos.~10--16, 19--24, 31--47, and 48--60, on Paschal encounter, the Church as celebrating subject, liturgical and symbolic formation, and ars celebrandi. \href{https://www.vatican.va/content/francesco/en/apost_letters/documents/20220629-lettera-ap-desiderio-desideravi.html}{Holy See English text}.
\end{itemize}

\subsection{Scripture, creeds, and principal patristic witnesses}

\begin{itemize}
  \item Principal biblical loci: Genesis 1--4, 14, and 22; Exodus 12--14, 16, 19--24, and 32--34; Leviticus 16; Deuteronomy 6, 8, and 10; Ruth 2; 2~Samuel 6; 1~Kings 8; 2~Maccabees 12:38--46; Nehemiah 8; Psalms 8, 22--23, 34, 51, 78, 85, 95, 104, 116, 118, 122, 136, 141, and 148; Isaiah 2, 6, 25--26, 52--55, and 61; Jeremiah 31; Lamentations 3; Baruch 3; Ezekiel 36--37; Daniel 3 and 7; Hosea 14; Malachi 1:11; Matthew 5--8, 15, 18, 20--21, and 25--28; Mark 8--11, 14, and 16; Luke 1--2, 4, 7, 18, 22, and 24; John 1, 6, 12--17, and 19--21; Acts 2, 4, 10, 12, and 20; Romans 5--6, 8, 10--12, and 14; 1~Corinthians 10--16; 2~Corinthians 5 and 13; Galatians 3; Ephesians 1--6; Philippians 2--3; Colossians 1 and 3--4; 1~Thessalonians 5; 1~Timothy 2; 2~Timothy 4; Hebrews 4 and 7--13; James 5; 1~Peter 2--3; 1~John 1; Revelation 4--5, 7--8, 19, and 21--22. Liturgical use is described as direct, typological, accommodative, or allusive according to context, not treated as automatic proof of textual ancestry.
  \item Councils of Nicaea I (325), Constantinople I (381), and Chalcedon (451): DS 125 and 150; ACO II.1.2, p.~128, for the received expanded creed as the faith of the 150 Fathers. Third Council of Toledo (589), canon~2, Mansi IX, 993, witnesses Western Eucharistic recitation; Berno of Reichenau, \emph{De quibusdam rebus ad missae officium pertinentibus}, PL 142:1060--1061, records Roman adoption under Benedict VIII in 1014.
  \item \emph{Didache}, chs.~8--10 and 14, for the Lord's Prayer, debated early meal thanksgivings, unity imagery, confession, and reconciliation before the Lord's Day breaking of bread. \href{https://www.newadvent.org/fathers/0714.htm}{Public English text}. Its meal prayers are not asserted to be the direct ancestor of a present anaphora.
  \item Justin Martyr, \emph{First Apology}, chs.~61 and 65--67, for baptismal Trinitarian practice, common intercession, kiss, mixed cup, presidential thanksgiving, Amen, distribution, Sunday readings, exhortation, gifts, and organized charity; \emph{Dialogue with Trypho} 41, for Malachi 1:11 and Christian thanksgiving. \href{https://www.newadvent.org/fathers/0126.htm}{\emph{First Apology}}. Justin supplies a second-century structural and apologetic witness, not the contemporary rubrics.
  \item Ignatius of Antioch, \emph{Letter to the Magnesians} 7, \emph{Letter to the Philadelphians} 4, \emph{Letter to the Smyrnaeans} 8, and \emph{Letter to the Trallians} 2.3, for one prayer, Eucharist, cup, altar, Bishop, presbytery, Deacons, and ministerial service in ecclesial unity. \href{https://www.newadvent.org/fathers/0105.htm}{Magnesians}; \href{https://www.newadvent.org/fathers/0108.htm}{Philadelphians}.
  \item Irenaeus of Lyons, \emph{Against Heresies} III.3.1--4, III.11.8, III.18--22, IV.14--18, and V.2, on apostolic proclamation, fourfold Gospel, recapitulation, divine pedagogy, created gifts, thanksgiving, purity of offering and offerer, ecclesial sacrifice, and bodily resurrection. \href{https://www.newadvent.org/fathers/0103318.htm}{Book III.18}; \href{https://www.newadvent.org/fathers/0103417.htm}{Book IV.17}; \href{https://www.newadvent.org/fathers/0103418.htm}{Book IV.18}.
  \item Tertullian, \emph{De corona} 3, for habitual Christian signing and care for the Eucharistic elements; \emph{De oratione} 11--18 and 27--28, for reconciliation, washing, posture, the kiss as seal of prayer, psalm response, and spiritual sacrifice; \emph{Adversus Marcionem} IV.9, CCSL 1, p.~560, for the Kyrie-related Gospel cry; and \emph{De praescriptione haereticorum} 13, for the received rule of faith. \href{https://www.newadvent.org/fathers/0304.htm}{\emph{De corona}}; \href{https://www.newadvent.org/fathers/0322.htm}{\emph{De oratione}}.
  \item Cyprian of Carthage, \emph{De dominica oratione}, especially 8--9, 18, and 31, on common filial prayer and \latin{Sursum corda}; Epistle 63 (numbered 62 in the linked collection), on the mixed cup, institution, and ecclesial unity; and \emph{De unitate ecclesiae} 5--7, on visible communion. \href{https://www.newadvent.org/fathers/050704.htm}{Lord's Prayer}; \href{https://www.newadvent.org/fathers/050662.htm}{mixed-cup letter}; \href{https://www.newadvent.org/fathers/050701.htm}{\emph{De unitate}}.
  \item The redactionally complex church order conventionally called the \emph{Apostolic Tradition}, especially ch.~4, for an anaphoral dialogue, Christological thanksgiving, institution, anamnesis and offering, Spirit petition, communion-unity petition, and doxology. Bradshaw, Johnson, and Phillips below control the cautions against assigning the reconstructed text to a single early-third-century Roman author.
  \item \emph{Apostolic Constitutions} VII.47 and VIII.5--15, a late-fourth-century Syrian pseudapostolic compilation, for Gloria-related hymnody, greeting, intercessions, kiss, gifts, Sanctus, epiclesis, Communion Amen and Psalm, thanksgiving, blessing, and dismissal. \href{https://www.newadvent.org/fathers/07158.htm}{Public English text of Book VIII}.
  \item The fourth-century Jerusalem catecheses, traditionally attributed to Cyril: \emph{Mystagogical Catechesis} II.4--6 and V.2--23, for baptismal signing, anointing, handwashing, peace, Preface dialogue, Sanctus, epiclesis, intercessions, Lord's Prayer, Communion chant, reception, Amen, and thanksgiving; \emph{Catechetical Lecture} 5.12, for the Creed as received ecclesial synthesis. \href{https://www.newadvent.org/fathers/310123.htm}{Mystagogical V}. These illuminate Eastern local orders and are not treated as Roman genetic sources.
  \item The late-fourth-century Milanese \emph{De sacramentis}, IV.4--6, especially IV.4.14--17, IV.5.21--24, and IV.6.26--27, for Christ's effective words and close Western parallels from the \latin{Quam oblationem} region through institution, anamnesis, offering, patriarchal sacrifices, and heavenly altar. The parallels do not date every earlier clause of the Roman Canon.
  \item Augustine of Hippo, Sermons 227, 256.1--3, 272, and 293.3; \emph{Enarrationes in Psalmos} 85.1 and 148.1; \emph{Tractatus in Ioannem} 1.8--9; \emph{Confessions} IX.11--13; and \emph{City of God} X.6, X.20, XIX.13, and XIX.17, for thanksgiving, sanctification, Psalter and Alleluia, Gospel witness, the Lord's Prayer, peace, Communion, the ecclesial Body and Amen, prayer for the dead, ecclesial sacrifice, ordered peace, and pilgrim use of earthly goods. \href{https://earlychurchtexts.com/public/augustine_sermon_272_eucharist.htm}{Sermon 272}; \href{https://www.newadvent.org/fathers/110109.htm}{\emph{Confessions} IX}; \href{https://www.newadvent.org/fathers/120119.htm}{\emph{City of God} XIX}.
  \item John Chrysostom, \emph{Homilies on First Corinthians} 24, on fraction, communion, one bread, and one body; \emph{Homilies on Matthew} 50.3--4 and 82, on Christ in the poor, institution, thanksgiving, Passion, worthy reception, and care for the needy. \href{https://www.newadvent.org/fathers/220124.htm}{First Corinthians 24}; \href{https://www.newadvent.org/fathers/200150.htm}{Matthew 50}; \href{https://www.newadvent.org/fathers/200182.htm}{Matthew 82}.
  \item Basil of Caesarea, \emph{De Spiritu Sancto} 27.66, for the baptismal orientation and the doxological confession of the Spirit; \emph{Homily Delivered in Time of Famine and Drought}, especially 6--8, PG 31:321--328, for the common claim of created goods and judgment of hoarding. \href{https://www.newadvent.org/fathers/3203.htm}{\emph{De Spiritu Sancto}}.
  \item Gregory Nazianzen, \emph{Oration} 38.17, on festival praise and Trinitarian glorification; Athanasius, \emph{Letter to Marcellinus} 10--12, on the Psalter as the worshipper's own voice; Origen, \emph{Disputatio cum Heracleida} 4.24, SC 67, p.~62, on the Kyrie-related acclamation; and \emph{Statuta Concilii Hipponensis breviata} 21, CCSL 149, p.~39, on the placement of the Gloria.
  \item \emph{First Clement} 59.2--61.3, for extensive common intercession ordered toward rulers and peace; Egeria, \emph{Itinerarium} 24.2, 6--7, 11 and 25.1--2, ed. Wilhelm Heraeus, \emph{Silviae vel potius Aetheriae peregrinatio ad loca sancta} (Heidelberg, 1908), for Jerusalem's announcements, blessings, dismissal, and congregation after the Eucharist.
  \item Gregory the Great, \emph{Registrum epistolarum} IX.26 (October 598), ed. Ewald--Hartmann, MGH \emph{Epistolae} II, pp.~59--60, for bounded Roman evidence concerning Alleluia, Kyrie and Christe, and placement of the Lord's Prayer; \emph{Regula pastoralis} III, prologue, for differentiated preaching to distinct hearers. \href{https://www.newadvent.org/fathers/360209012.htm}{Public English arrangement of the letter as IX.12}. Gregory describes and regulates these practices; he is not called their composer.
  \item Thomas Aquinas, \emph{Summa theologiae} III, q.~78, aa.~1--4; q.~82, a.~1; q.~83, a.~3, ad~2; and q.~83, a.~4 and ad~8, for Christ's effective words, ministerial presidency, altar symbolism, the Canon's ordered unity, and the accepted biblical offerings. \href{https://www.newadvent.org/summa/4078.htm}{Question 78}; \href{https://www.newadvent.org/summa/4082.htm}{Question 82}; \href{https://www.newadvent.org/summa/4083.htm}{Question 83}.
\end{itemize}

\subsection{Roman development and reform records}

\begin{itemize}[nosep]
  \item \emph{Ordo Romanus I}, in Michel Andrieu, \emph{Les Ordines Romani du haut moyen âge}, vol.~II, \emph{Les textes (Ordines I--XIII)} (Louvain, 1948), pp.~67--108, especially nn.~8--21, for the layered late-seventh-/early-eighth-century papal stational order from entrance through dismissal. It is not treated as a parish order, a single-date composition, or ``the Mass of Gregory.''
  \item The sacramentary traditions conventionally called Veronese, Old Gelasian, and Gregorian, in L.~C. Mohlberg et al., eds., \emph{Sacramentarium Veronense (Cod. Bibl. Capit. Veron. LXXXV [80])} (Rome, 1956); Mohlberg et al., eds., \emph{Liber sacramentorum Romanae aeclesiae ordinis anni circuli (Cod. Vat. Reg.~lat. 316)} (Rome, 1960); and Jean Deshusses, ed., \emph{Le Sacramentaire grégorien}, 3 vols. (Fribourg, 1971--1982). They establish late-antique and early-medieval Roman corpora of Collects, prayers over gifts, Prefaces, post-Communion prayers, and prayers over the people. Manuscript date is distinguished from the possible age of transmitted prayer material.
  \item Codex Alexandrinus, British Library Royal MS 1 D V--VIII, and the Bangor Antiphonary, in F.~E. Warren, ed., \emph{The Antiphonary of Bangor}, Henry Bradshaw Society 4 and 10 (1893--1895), for Gloria witnesses; Notker Balbulus, preface to \emph{Liber hymnorum}, for one account of textual aids to long Alleluia melodies. None supplies a universal origin story for the present Roman Gloria or every Sequence.
  \item \emph{Liber pontificalis}, ed. Louis Duchesne, vol.~I, for retrospective notices concerning Celestine, early papal Gloria and Sanctus use, and Sergius I's Agnus Dei. These notices are reception evidence, not contemporary autograph decrees.
  \item Amalarius of Metz, \emph{Liber officialis} III.19.26, ed. J.-M. Hanssens, vol.~II, p.~319, for a ninth-century witness to offertory incensation described as not then Roman; the source checks a claim of universal primitive Roman use.
  \item The Christmas-collect/mixing-prayer genealogy in Ludovico Antonio Muratori, \emph{Liturgia Romana Vetus}, vol.~I, p.~467; Edmond Martène, \emph{De antiquis Ecclesiae ritibus}, vol.~I, pp.~510--511; and Jungmann, vol.~II, pp.~57--64; Western ninth- to eleventh-century witnesses to \latin{In spiritu humilitatis}; and Robert Lippe, ed., \emph{Missale Romanum Mediolani, 1474}, Henry Bradshaw Society 17 and 33 (London, 1899/1907), pp.~198--212, for the received \latin{Orate, fratres} pair. These establish bounded formula histories, not one apostolic Offertory.
  \item Victor Leroquais, ``L'Ordo Missae du sacramentaire d'Amiens,'' \emph{Ephemerides Liturgicae} 41 (1927), p.~444; Leroquais, \emph{Les sacramentaires et les missels manuscrits des bibliothèques publiques de France}, vol.~I, p.~30; and Richter--Schönfelder, eds., \emph{Sacramentarium Fuldense}, no.~24, for medieval private Communion preparations. The current Missal retains two approved prayers from a much larger medieval family.
  \item Provincial Council of Aix (September 1585), Eucharistic constitution, Labbe--Cossart XV, cols.~1119ff., as a terminus for the faithful's showing formula; Manlio Sodi and Juan Javier Flores Arcas, eds., \emph{Rituale Romanum: Editio princeps (1614)} (Vatican City, 2004), tit.~IV, cap.~1, for the received \latin{Ecce Agnus Dei} and threefold \latin{Domine, non sum dignus} complex. No unverified internal column is invented for Aix.
  \item Council of Trent, Pius V's apostolic constitution \emph{Quo primum} (14 July 1570), and the 1570 \emph{Missale Romanum}, read as doctrinal definition and authoritative regularization of inherited Roman/curial worship rather than composition of all its texts or prohibition of every later reform.
  \item Pius X, motu proprio \emph{Tra le sollecitudini} (22 November 1903), introduction and no.~3, on participation and restoration of congregational Gregorian chant, \href{https://www.vatican.va/content/pius-x/it/motu_proprio/documents/hf_p-x_motu-proprio_19031122_sollecitudini.html}{Holy See Italian text}; Sacred Congregation of the Council, decree \emph{Sacra Tridentina Synodus} (20 December 1905), especially norms 1--6 on frequent Communion, ASS 38 (1905--1906), 400--406, \href{https://www.vatican.va/archive/ass/documents/ASS-38-1905-6-ocr.pdf}{official scan}.
  \item Pius XII, apostolic constitution \emph{Christus Dominus} (6 January 1953), AAS 45 (1953), 15--24, \href{https://www.vatican.va/content/pius-xii/la/apost_constitutions/documents/hf_p-xii_apc_19530106_christus-dominus.html}{Holy See Latin text}; motu proprio \emph{Sacram Communionem} (19 March 1957), AAS 49 (1957), 177--178, for Eucharistic-fast and evening-Mass discipline; Sacred Congregation of Rites, decree \emph{Dominicae Resurrectionis} (9 February 1951), AAS 43 (1951), 128--137; decree \emph{Maxima redemptionis nostrae mysteria} (16 November 1955), AAS 47 (1955), 838--847. These acts establish immediate institutional horizons for twentieth-century reform before Vatican II.
  \item Sacred Congregation of Rites and Consilium, instruction \emph{Inter Oecumenici} (26 September 1964), AAS 56 (1964), 877--900; instruction \emph{Tres abhinc annos} (4 May 1967), AAS 59 (1967), 442--448; the 1965, 1967, and 1969 transitional Orders; and the 1970 first typical edition, used as reform chronology rather than silently conflated with the current 2008 text.
  \item Sacred Congregation of Rites, decree \emph{Novis hisce temporibus} (13 November 1962), AAS 54 (1962), 873, inserting Saint Joseph into the Roman Canon; the later \emph{Paternas vices} decree controls the distinct 2013 insertion into Prayers II--IV.
  \item Sacred Congregation of Rites, decree and norms for Eucharistic Prayers I--IV, \emph{Notitiae} 4 (1968), pp.~156--160, \href{https://www.cultodivino.va/content/dam/cultodivino/rivista-notitiae/1960/notitiae-04-\%281968\%29/Notitiae-040-1968.pdf}{official journal scan}.
  \item Anton Hänggi and Irmgard Pahl, eds., \emph{Prex Eucharistica}, vol.~I, 3rd ed., rev. Heinzgerd Brakmann and Albert Gerhards (Fribourg, 1998), pp.~347--357, for the Alexandrian Basil anaphora used as a comparative source in the composition history of Eucharistic Prayer IV; comparison is not treated as direct translation.
  \item Sacred Congregation for Divine Worship, \emph{Postquam de precibus} and the 1975 Reconciliation and children's-prayer dossier, \emph{Notitiae} 11 (1975), pp.~4--12, \href{https://www.cultodivino.va/content/dam/cultodivino/rivista-notitiae/1970/notitiae-11-\%281975\%29/Notitiae-101-1975.pdf}{official journal scan}; Congregation for the Sacraments and Divine Worship, 1983 faculty and text for Reconciliation I--II, \emph{Notitiae} 19 (1983), pp.~270--279, \href{https://www.cultodivino.va/content/dam/cultodivino/rivista-notitiae/1980/notitiae-19-\%281983\%29/Notitiae-202-1983.pdf}{official journal scan}.
  \item Congregation for Divine Worship and the Discipline of the Sacraments, normative Latin recension and \emph{Praenotanda} for the Eucharistic Prayer for Various Needs, \emph{Notitiae} 27 (1991), pp.~388--459, including Pere Tena's commentary and Anton Hänggi's history, \href{https://www.cultodivino.va/content/dam/cultodivino/rivista-notitiae/1990/notitiae-27-\%281991\%29/Notitiae-301-1991.pdf}{official journal scan}.
  \item USCCB Committee on Divine Worship, \href{https://www.usccb.org/resources/newsletter-2012-02.pdf}{Newsletter, February 2012}, pp.~6--8, for a calibrated composition account of Eucharistic Prayers II--IV; \href{https://www.usccb.org/resources/newsletter-2014-02-and-03.pdf}{Newsletter, February--March 2014}, for the Preface discipline of the Reconciliation prayers.
\end{itemize}

\subsection{Modern historical and theological studies}

\begin{itemize}[nosep]
  \item Josef A. Jungmann, \emph{The Mass of the Roman Rite: Its Origins and Development}, trans. Francis A. Brunner, 2 vols. (New York: Benziger, 1951--1955). Used critically for diachronic Roman synthesis beside primary witnesses; claims tied to later-discovered sources or reform history are not adopted uncritically.
  \item Paul F. Bradshaw, Maxwell E. Johnson, and L. Edward Phillips, \emph{The Apostolic Tradition: A Commentary} (Minneapolis: Fortress Press, 2002), for text, reconstruction, redactional complexity, and cautions concerning Hippolytan authorship, Roman provenance, and date.
  \item Enrico Mazza, \emph{The Eucharistic Prayers of the Roman Rite}, trans. Matthew J. O'Connell (New York: Pueblo, 1986), for anaphoral structure and history, checked against the received books and official reform dossiers.
  \item Annibale Bugnini, \emph{The Reform of the Liturgy, 1948--1975}, trans. Matthew J. O'Connell (Collegeville: Liturgical Press, 1990), used as an involved reformer's documentary narrative rather than a neutral or solitary authority.
  \item Cipriano Vagaggini, \emph{The Canon of the Mass and Liturgical Reform}, trans. Peter Coughlan (Staten Island: Alba House, 1967), for the theological and compositional horizon of the new Eucharistic Prayers, read with the official acts and later historical cautions.
  \item Robert F. Taft, \emph{Beyond East and West: Problems in Liturgical Understanding}, 2nd ed. (Rome: Pontifical Oriental Institute, 1997), for comparative method that respects local traditions and resists using Eastern structures as unqualified Roman genealogy.
\end{itemize}

\endgroup
